\chapter{Landasan Teori}

\section{Tinjauan Pustaka}
Penelitian terkait prediksi Diabetes Mellitus Tipe 2 (T2DM) menggunakan \eng{machine learning} telah banyak dilakukan dengan berbagai pendekatan. Secara umum, penelitian-penelitian tersebut berfokus pada peningkatan akurasi model, pemilihan algoritma yang tepat, serta bagaimana model dapat digunakan secara lebih efektif dalam konteks kesehatan.

Beberapa penelitian awal fokus pada perbandingan algoritma \eng{machine learning}. Penelitian \cite{ref10} membandingkan \eng{Logistic Regression}, \eng{Random Forest}, \eng{Support Vector Machine} (SVM), dan \eng{K-Nearest Neighbors} (KNN) dengan menerapkan SMOTE untuk mengatasi ketidakseimbangan data. Hasilnya menunjukkan bahwa \eng{Random Forest} memiliki performa terbaik dengan \eng{F1-score} sebesar 96,92\% dan AUC 99,70\%. Meskipun SVM menunjukkan performa yang baik, waktu komputasi yang lebih besar menjadi kendala. Temuan ini mengindikasikan bahwa metode \eng{ensemble} lebih unggul dibandingkan dengan algoritma tunggal. Hal ini juga diperkuat oleh penelitian \cite{ref14} yang membandingkan beberapa algoritma seperti \eng{Logistic Regression}, \eng{Decision Tree}, \eng{Random Forest}, \eng{Gradient Boosting}, XGBoost, dan LightGBM. Hasil penelitian menunjukkan bahwa LightGBM memberikan akurasi tertinggi sebesar 95,20\%. Temuan ini semakin memperkuat bahwa algoritma berbasis \eng{boosting} mampu menangani data yang kompleks dengan lebih baik.

Selain pemilihan algoritma, optimasi model juga berperan penting dalam meningkatkan performa. Penelitian \cite{ref22} menunjukkan bahwa penggunaan \eng{feature engineering} dan \eng{hyperparameter tuning} dapat meningkatkan performa model secara signifikan. Dengan menggunakan RFECV dan GridSearchCV, model XGBoost berhasil mencapai akurasi sebesar 94\%, lebih tinggi dibandingkan model lainnya. Hal ini menunjukkan bahwa kualitas fitur dan parameter sangat berpengaruh terhadap hasil prediksi. Penggunaan metode \eng{ensemble} juga banyak diteliti karena kemampuannya dalam meningkatkan performa model. Penelitian \cite{ref26} menunjukkan bahwa metode \eng{ensemble} seperti \eng{Bagging}, \eng{Boosting}, dan \eng{Stacking} mampu mencapai akurasi hingga 97\%, lebih tinggi dibandingkan model individual. Hal ini terjadi karena \eng{ensemble} menggabungkan kelebihan dari beberapa model sehingga menghasilkan prediksi yang lebih stabil.

Penelitian lain \cite{ref27} mengombinasikan \eng{feature selection} dengan algoritma \eng{boosting} seperti LightGBM. Hasilnya menunjukkan akurasi sebesar 85,16\% dan \eng{F1-score} sebesar 85,41\%, dengan jumlah fitur yang lebih sedikit serta waktu pelatihan yang lebih efisien. Hal ini menunjukkan bahwa pemilihan fitur yang tepat dapat meningkatkan efisiensi tanpa mengurangi performa model. Dalam hal optimasi model, penelitian \cite{ref28} mengembangkan metode berbasis \eng{ensemble} yang mampu mencapai akurasi hingga 99\%, dengan \eng{precision} 98,5\% dan \eng{recall} 99\%. Hasil ini menunjukkan bahwa kombinasi antara optimasi parameter dan \eng{ensemble learning} dapat menghasilkan model dengan performa yang sangat tinggi.

Selain itu, permasalahan ketidakseimbangan data juga menjadi perhatian dalam beberapa penelitian. Penelitian \cite{ref29} mengombinasikan SMOTE dan Near Miss serta menerapkan \eng{Explainable AI}. Hasilnya menunjukkan bahwa model XGBoost dengan SMOTE mampu mencapai akurasi 99\% dan \eng{F1-score} 1,00. Penelitian \cite{ref16} juga menunjukkan bahwa penggunaan SMOTE dapat meningkatkan performa model, dengan nilai AUC mencapai 0,89 dan akurasi sebesar 90,2\%. Hal ini menunjukkan bahwa \eng{preprocessing} data memiliki peran penting dalam meningkatkan kualitas model.

Di sisi lain, interpretabilitas model juga menjadi aspek penting, terutama dalam bidang kesehatan. Penelitian \cite{ref30} menunjukkan bahwa pendekatan \eng{interpretable machine learning} dapat membantu mengidentifikasi faktor-faktor penting dalam diagnosis diabetes. Sementara itu, penelitian \cite{ref15} menggabungkan AutoML dan \eng{Explainable AI} untuk menghasilkan model yang tidak hanya akurat tetapi juga transparan.

Secara keseluruhan, berbagai penelitian ini menunjukkan bahwa \eng{machine learning} memiliki potensi besar dalam prediksi diabetes. Namun, sebagian besar penelitian masih berfokus pada satu aspek tertentu, seperti pemilihan algoritma, optimasi model, atau interpretabilitas. Selain itu, perbandingan antara algoritma \eng{baseline} dan \eng{boosting} dalam satu kerangka yang sama masih terbatas. Oleh karena itu, penelitian ini bertujuan untuk menggabungkan pendekatan-pendekatan tersebut, yaitu perbandingan algoritma, \eng{hyperparameter tuning} menggunakan Optuna, serta penerapan \eng{Explainable AI}, untuk menghasilkan model yang lebih akurat dan mudah dipahami.

\begin{landscape}
\scriptsize
\setlength{\tabcolsep}{3pt}
\begin{longtable}{|C{0.7cm}|L{2.9cm}|L{5.4cm}|L{3.0cm}|L{4.4cm}|L{3.7cm}|}
\caption{Ringkasan Penelitian Terdahulu}\label{tab:penelitian-terdahulu}\\
\hline
\rowcolor{gray!25}\textbf{No.} & \textbf{Judul dan Tahun} & \textbf{Pembahasan} & \textbf{Metode} & \textbf{Kelebihan dan Kekurangan} & \textbf{Perbedaan Penelitian}\\ \hline
\endfirsthead
\multicolumn{6}{c}{\tablename\ \thetable\ -- lanjutan dari halaman sebelumnya}\\
\hline
\rowcolor{gray!25}\textbf{No.} & \textbf{Judul dan Tahun} & \textbf{Pembahasan} & \textbf{Metode} & \textbf{Kelebihan dan Kekurangan} & \textbf{Perbedaan Penelitian}\\ \hline
\endhead
\hline\multicolumn{6}{r}{Bersambung ke halaman berikutnya}\\
\endfoot
\hline
\endlastfoot

1 & \eng{Comparison of Logistic Regression, Random Forest, Support Vector Machine (SVM) and K-Nearest Neighbor (KNN) Algorithms in Diabetes Prediction}, 2025 \cite{ref10} & Penelitian ini menunjukkan bahwa \eng{Random Forest} mampu mencapai \eng{F1-Score} sebesar 96,92\% dengan AUC 99,70\% pada dataset berukuran 100 ribu data, sehingga melampaui performa \eng{Logistic Regression} sekitar 75\%, SVM 74,4\%, dan KNN 70,8\%. Penerapan SMOTE berhasil menyeimbangkan data dengan meningkatkan jumlah sampel diabetes dari 15\% menjadi distribusi seimbang 85K:85K. Hasil ini memberikan bukti empiris bahwa metode \eng{ensemble} seperti \eng{Random Forest} lebih unggul dibandingkan algoritma sederhana. & \eng{Logistic Regression}, \eng{Random Forest}, SVM, KNN, 80:20 \eng{train-test split}, 5-fold \eng{stratified cross-validation}, SMOTE untuk \eng{class balancing}. & \textbf{Kelebihan:} Evaluasi sistematis empat algoritma klasik pada dataset besar; efektivitas SMOTE untuk \eng{class imbalance} terbukti empiris; konsistensi RF tinggi. \newline\textbf{Kekurangan:} Tanpa \eng{hyperparameter tuning} sistematis, hanya parameter bawaan, tanpa analisis XAI/interpretabilitas, tanpa analisis \eng{feature importance}, dan tidak ada validasi lintas dataset. & Penelitian ini tidak hanya membandingkan algoritma \eng{baseline}, tetapi juga melakukan perbandingan eksplisit antara algoritma \eng{baseline} dan algoritma berbasis \eng{boosting} dengan total enam algoritma. Selain itu, penelitian ini menerapkan \eng{hyperparameter tuning} menggunakan Optuna serta mengintegrasikan SHAP.\\ \hline

2 & \eng{Diabetes Prediction Using Machine Learning with Feature Engineering and Hyperparameter Tuning}, 2024 \cite{ref22} & Kombinasi XGBoost dengan RFECV mencapai akurasi 94\%. RFECV mengidentifikasi lima fitur utama, yaitu \eng{Glucose}, BMI, \eng{Age}, DPF, dan \eng{Skin Thickness}. Proses \eng{feature engineering} meningkatkan performa 2--3\% dibandingkan \eng{baseline}. & Lima teknik FSA (SelectKBest, RFE, RFECV, \eng{Forward/Backward}), lima algoritma \eng{boosting} (RF, GB, XGB, LGBM, CatBoost), GridSearchCV, dataset PIDD. & \textbf{Kelebihan:} Perbandingan \eng{feature selection} komprehensif, optimasi parameter terperinci, dan panduan praktis. \newline\textbf{Kekurangan:} Hanya satu dataset PIMA, representasi demografis terbatas, tanpa XAI, tanpa perbandingan algoritma \eng{baseline}, dan akurasi masih berpotensi ditingkatkan. & Penelitian ini menambahkan pendekatan optimasi melalui Optuna yang lebih efisien, mencakup perbandingan algoritma \eng{baseline}, dan mengintegrasikan SHAP untuk memahami hubungan antar fitur.\\ \hline

3 & Evaluasi Model \eng{Ensemble Learning} pada Identifikasi Faktor Risiko Diabetes Mellitus, 2025 \cite{ref26} & \eng{Gradient Boosting} mencapai nilai $R^2$ sebesar 0,65 dan seluruh metode \eng{ensemble} mencapai akurasi sekitar 97\%. Naïve Bayes cenderung menghasilkan prediksi berlebih dengan \eng{recall} tinggi namun \eng{precision} rendah, sedangkan \eng{Logistic Regression} dan SVM lebih konservatif. & Lima \eng{base learners} (NB, LR, SVM, DT, KNN), empat metode \eng{ensemble} (Bagging, Boosting, Stacking, Voting), KNN Imputer, SMOTE, dataset PIMA. & \textbf{Kelebihan:} Perbandingan \eng{ensemble} komprehensif dan analisis \eng{sensitivity-specificity trade-off}. \newline\textbf{Kekurangan:} Hanya satu dataset, demografi terbatas, tanpa rincian \eng{feature importance}, integrasi interpretabilitas klinis minimal, dan tidak ada perbandingan eksplisit \eng{baseline-boosting}. & Penelitian ini menambahkan perbandingan eksplisit antara algoritma \eng{baseline} dan \eng{boosting}, memperluas evaluasi, serta menggunakan SHAP untuk memahami kontribusi fitur.\\ \hline

4 & \eng{Diabetes Prediction Using Feature Selection Algorithms and Boosting-Based Machine Learning Classifiers}, 2025 \cite{ref27} & LightGBM yang dikombinasikan dengan Boruta mencapai akurasi 85,16\%, \eng{F1-score} 85,41\%, dan AUC 0,9052 pada PIDD. Waktu pelatihan turun 54,96\% dengan hanya lima dari delapan fitur. Validasi eksternal pada DiaHealth menghasilkan akurasi 99,39\% dan \eng{recall} 100\%. & Lima FSA (RFE, GWO, PSO, GA, Boruta), dua \eng{boosting} (XGB, LGBM), 5-fold CV, 70:30 split, SHAP, PIDD dan DiaHealth. & \textbf{Kelebihan:} Integrasi FSA dan XAI, validasi eksternal kuat, efisiensi komputasi meningkat. \newline\textbf{Kekurangan:} Terbatas pada dua dataset, kedalaman SHAP minimal, tidak ada perbandingan \eng{baseline}, dan tidak ada interpretasi lokal tingkat pasien. & Penelitian ini tetap menggunakan SHAP, namun mengombinasikannya dengan perbandingan algoritma \eng{baseline} dan \eng{boosting}, evaluasi lebih sistematis, dan pengembangan model yang lebih aplikatif.\\ \hline

5 & \eng{An Optimized Diabetes Mellitus Detection Model for Improved Prediction of Accuracy and Clinical Decision-Making}, 2024 \cite{ref28} & Model GBM-DRU menghasilkan akurasi 99\%, \eng{precision} 98,5\%, dan \eng{recall} 99\%, lebih unggul dibandingkan RF, J48, SVM, XGBoost, Naïve Bayes, dan \eng{Soft Voting}. Hasil \eng{cross-validation} menunjukkan stabilitas 95\%, tetapi transparansi model dan interpretasi klinis belum tersedia. & \eng{Gradient Boosting} + \eng{Data Reduction Unit}, \eng{feature engineering}, \eng{ensemble stacking}, 10-fold CV, spesifikasi dataset tidak terdokumentasi. & \textbf{Kelebihan:} Akurasi sangat tinggi dan hasil validasi silang kuat. \newline\textbf{Kekurangan:} Reprodusibilitas dataset menjadi perhatian, tanpa XAI, potensi \eng{overfitting}, generalisasi tidak jelas, dan tanpa interpretasi fitur. & Penelitian ini menekankan keseimbangan antara performa dan transparansi, menggunakan optimasi \eng{hyperparameter} serta evaluasi yang lebih sistematis untuk meminimalkan \eng{overfitting}.\\ \hline

6 & \eng{SMOTE Oversampling and Near Miss Undersampling Based Diabetes Diagnosis from Imbalanced Dataset with XAI Visualization}, 2023 \cite{ref29} & XGBoost dengan SMOTE mencapai akurasi 99\% dan \eng{F1-score} 1,00. Integrasi SHAP dan LIME meningkatkan kepercayaan pengguna serta keandalan model dalam konteks klinis. & Sepuluh model ML, SMOTE, Near Miss, teknik XAI (SHAP dan LIME), dataset diabetes tidak seimbang. & \textbf{Kelebihan:} \eng{F1-score} sempurna, adopsi XAI, dan mitigasi \eng{class imbalance}. \newline\textbf{Kekurangan:} Perbandingan algoritma terbatas, optimasi parameter minimal, pengujian generalisasi tidak jelas, tanpa validasi multi-dataset, dan interaksi fitur tidak dieksplorasi. & Penelitian ini mengintegrasikan optimasi model menggunakan Optuna serta menerapkan SHAP sebagai metode interpretasi utama, dengan evaluasi yang lebih sistematis.\\ \hline

7 & \eng{Predictive Analysis of Diabetes-Risk with Class Imbalance}, 2022 \cite{ref16} & Model ANN dengan SMOTE mencapai akurasi 90,2\% dan AUC 0,89. SMOTE efektif menangani ketidakseimbangan data, sedangkan Tomek-links memberikan pengaruh minimal. Resistensi insulin dan BMI tinggi ditemukan sebagai faktor risiko utama. & SMOTE, Tomek-links, ANN, pemeringkatan fitur (odds ratio, Boruta, varImp), \eng{cost-sensitive learning}, dataset khusus perempuan. & \textbf{Kelebihan:} Membandingkan berbagai strategi ketidakseimbangan dan menghasilkan wawasan klinis. \newline\textbf{Kekurangan:} Hanya satu gender, dataset kecil, tanpa perbandingan sistematis \eng{baseline-boosting}, tanpa eksplorasi interaksi fitur, dan generalisasi terbatas. & Penelitian ini memperluas pendekatan pada perbandingan algoritma \eng{baseline} dan \eng{boosting}, serta mengintegrasikan interpretasi model untuk memahami faktor risiko diabetes.\\ \hline

8 & \eng{Prediction of Type-2 Diabetes Mellitus Disease Using Machine Learning Classifiers and Techniques}, 2022 \cite{ref14} & LightGBM mencapai akurasi tertinggi 95,20\%, diikuti Extra Trees 94,60\%, Random Forest 94,80\%, Decision Tree 94,40\%, Gradient Boosting 94,10\%, XGBoost 83,30\%, dan Logistic Regression 75,20\%. & Tujuh algoritma (LR, XGB, GB, DT, ET, RF, LGBM), dataset PIDD, 80:20 split, metrik evaluasi dasar. & \textbf{Kelebihan:} Perbandingan tujuh algoritma dan panduan pemilihan praktis. \newline\textbf{Kekurangan:} Tanpa \eng{hyperparameter tuning}, analisis fitur minimal, tanpa XAI, tanpa laporan \eng{cross-validation}, dan hanya satu dataset. & Penelitian ini menambahkan \eng{hyperparameter tuning} menggunakan Optuna serta mengintegrasikan interpretabilitas menggunakan SHAP.\\ \hline

9 & \eng{Unveiling Diagnostic Information for Type 2 Diabetes through Interpretable Machine Learning}, 2025 \cite{ref30} & Model DSpaLaRefiner mencapai akurasi 98,79\% pada SWMU, 97,22\% pada NSMC, 97,67\% pada data gabungan, dan validasi eksternal 91,90\%. Model mengidentifikasi kasus T2DM pada kadar glukosa normal dengan indikator sekunder. Fitur penting meliputi GLU, usia, glukosa urin, HCT, dan HDL-C. & DSpaLaRefiner, 39 fitur, dua dataset rumah sakit, ekstraksi informasi intrinsik, serta algoritma pembanding SVM, RF, MLP, dan AN. & \textbf{Kelebihan:} Akurasi dan interpretabilitas global-lokal tinggi, rentang fitur klinis teridentifikasi, validasi eksternal kuat. \newline\textbf{Kekurangan:} Hanya data tabular, terbatas pada dua rumah sakit, klasifikasi biner, dan kompleksitas algoritma tinggi. & Penelitian ini menggunakan pendekatan yang lebih sederhana melalui perbandingan algoritma \eng{machine learning}, dengan penekanan pada kemudahan interpretasi dan implementasi praktis.\\ \hline

10 & \eng{Towards Transparent Diabetes Prediction: Combining AutoML and Explainable AI for Improved Clinical Insights}, 2025 \cite{ref15} & AutoGluon menghasilkan rata-rata akurasi 85,01\% pada PIDD, Sklearn, dan Rural African-American. Analisis SHAP menunjukkan \eng{Glucose} dan BMI sebagai fitur paling berpengaruh. LIME digunakan untuk interpretasi individual, \eng{Counterfactual Analysis} untuk skenario hipotetis, dan aplikasi Streamlit memperluas akses. & AutoGluon, SHAP, LIME, IG, AM, CA, LightGBM, XGBoost, CatBoost, RF, 5-fold \eng{stratified CV}, tiga dataset, dan aplikasi Streamlit. & \textbf{Kelebihan:} Otomasi AutoML, lima teknik XAI, generalisasi tiga dataset, dan alat klinis interaktif. \newline\textbf{Kekurangan:} Spesifisitas demografi PIDD, akurasi moderat pada dataset utama, dan waktu pelatihan AutoML singkat dapat membatasi eksplorasi. & Penelitian ini lebih berfokus pada perbandingan algoritma \eng{baseline} dan \eng{boosting}, penggunaan Optuna untuk optimasi, dan SHAP sebagai metode interpretasi utama.\\ \hline
\end{longtable}
\setlength{\tabcolsep}{6pt}
\end{landscape}

\section{Landasan Teori}
Landasan teori ini menjelaskan konsep-konsep dasar yang digunakan dalam penelitian, meliputi Diabetes Mellitus Tipe 2, \eng{machine learning} dan klasifikasi, algoritma \eng{baseline} dan \eng{boosting}, \eng{hyperparameter tuning} dengan Optuna, evaluasi model, serta \eng{Explainable AI} menggunakan SHAP.

\subsection{Diabetes Mellitus Tipe 2}
Diabetes Mellitus Tipe 2 (T2DM) merupakan penyakit metabolik yang ditandai dengan tingginya kadar glukosa dalam darah akibat gangguan kerja insulin dalam tubuh \cite{ref6}. T2DM merupakan jenis diabetes yang paling umum, dengan proporsi sekitar 90--95\% dari seluruh kasus diabetes di dunia \cite{ref7}. Jumlah penderita T2DM terus meningkat dari tahun ke tahun dan menjadi salah satu masalah kesehatan global yang serius karena dapat menyebabkan berbagai komplikasi serta meningkatkan angka kematian \cite{ref7,ref8}. Secara medis, T2DM terjadi akibat adanya resistensi insulin dan penurunan fungsi sel beta pankreas, sehingga tubuh tidak mampu mengatur kadar glukosa darah secara optimal \cite{ref6}. Resistensi insulin menyebabkan sel-sel tubuh menjadi kurang responsif terhadap insulin, sedangkan disfungsi sel beta pankreas menyebabkan produksi insulin menurun secara bertahap. Kombinasi kedua kondisi ini mengakibatkan kadar gula darah tetap tinggi dalam jangka waktu yang lama.

Kondisi tersebut dapat memicu berbagai komplikasi serius seperti penyakit kardiovaskular, neuropati, gangguan ginjal, serta gangguan penglihatan \cite{ref8}. Selain itu, T2DM juga dipengaruhi oleh berbagai faktor risiko seperti faktor genetik, usia, obesitas, pola makan tidak sehat, dan kurangnya aktivitas fisik \cite{ref9}. Beberapa penelitian juga menunjukkan bahwa biomarker tertentu, seperti asam amino dan asam lemak, memiliki hubungan dengan peningkatan risiko T2DM, yang menandakan bahwa penyakit ini bersifat kompleks dan dipengaruhi oleh berbagai faktor biologis maupun gaya hidup \cite{ref9}. Oleh karena itu, deteksi dini dan penanganan yang tepat menjadi sangat penting untuk mencegah perkembangan penyakit serta mengurangi risiko komplikasi yang lebih lanjut.

\subsection{Machine Learning dan Klasifikasi}
\eng{Machine learning} merupakan bagian dari \eng{Artificial Intelligence} yang memungkinkan sistem komputer untuk belajar dari data dan meningkatkan kinerjanya tanpa harus diprogram secara eksplisit. Istilah \eng{machine learning} pertama kali diperkenalkan oleh Arthur Samuel pada tahun 1959, yang mendefinisikannya sebagai kemampuan komputer untuk belajar secara otomatis dari pengalaman. Seiring perkembangan teknologi, \eng{machine learning} telah banyak digunakan dalam berbagai bidang, termasuk kesehatan, karena kemampuannya dalam mengolah data dalam jumlah besar dan menemukan pola yang kompleks \cite{ref1}.

Dalam bidang kesehatan, \eng{machine learning} dimanfaatkan untuk membantu proses prediksi dan pengambilan keputusan klinis, terutama dengan memanfaatkan data rekam medis elektronik (EHR). Data EHR memungkinkan model untuk mempelajari kondisi pasien secara lebih komprehensif, sehingga dapat meningkatkan akurasi prediksi penyakit \cite{ref2,ref3}. Salah satu pendekatan yang umum digunakan adalah \eng{supervised learning}, yaitu metode pembelajaran di mana model dilatih menggunakan data yang telah memiliki label. Salah satu bentuk permasalahan dalam \eng{supervised learning} adalah klasifikasi, yaitu proses mengelompokkan data ke dalam kategori tertentu.

Dalam penelitian ini, \eng{machine learning} digunakan untuk klasifikasi biner, yaitu menentukan apakah seseorang termasuk dalam kategori penderita T2DM atau tidak. Algoritma \eng{baseline} seperti \eng{Logistic Regression}, SVM, dan KNN sering digunakan karena relatif sederhana \cite{ref10,ref11}. Sementara itu, teknik \eng{boosting} seperti \eng{Gradient Boosting}, XGBoost, dan LightGBM mampu meningkatkan performa dengan menggabungkan sejumlah model sederhana menjadi model yang lebih kuat \cite{ref16,ref20}.

\begin{figure}[H]
\centering
\includegraphics[width=0.9\textwidth]{perbandingan_algoritma.png}
\caption{Perbandingan Karakteristik Algoritma \eng{Machine Learning}}
\label{fig:perbandingan-algoritma}
\end{figure}

Gambar \ref{fig:perbandingan-algoritma} menunjukkan perbedaan karakteristik antara algoritma \eng{baseline} dan algoritma berbasis \eng{boosting}. Algoritma \eng{baseline} seperti \eng{Logistic Regression}, SVM, dan KNN umumnya lebih sederhana dan mudah diinterpretasikan, namun memiliki keterbatasan dalam menangani hubungan non-linear yang kompleks. Sebaliknya, algoritma \eng{boosting} seperti \eng{Gradient Boosting}, XGBoost, dan LightGBM mampu menangani data yang lebih kompleks melalui pendekatan \eng{ensemble}, sehingga cenderung menghasilkan performa prediksi yang lebih tinggi \cite{ref31}.

\subsection{Normalisasi Data}
Normalisasi data merupakan proses transformasi nilai pada fitur numerik ke dalam skala tertentu sehingga setiap variabel memiliki distribusi atau rentang nilai yang lebih seragam. Proses ini bertujuan untuk mengurangi perbedaan skala antar fitur, sehingga setiap variabel memiliki kontribusi yang seimbang dalam proses pembelajaran model \cite{ref32}. Perbedaan skala yang terlalu besar antar fitur dapat menyebabkan algoritma \eng{machine learning} menjadi bias terhadap fitur tertentu yang memiliki nilai dominan, sehingga berpotensi menurunkan performa model.

Normalisasi umumnya diterapkan pada algoritma yang sensitif terhadap skala data, seperti KNN, SVM, \eng{Neural Network}, serta metode berbasis \eng{clustering}. Salah satu metode yang digunakan dalam penelitian ini adalah \eng{Standard Scaler} atau standarisasi berbasis Z-score. Rumusnya dinyatakan sebagai:
\begin{equation}
 z = \frac{x-\mu}{\sigma},
 \label{eq:zscore}
\end{equation}
dengan $x$ menyatakan nilai data, $\mu$ sebagai rata-rata, $\sigma$ sebagai standar deviasi, dan $z$ merupakan hasil transformasi \cite{ref32}.

\subsection{Logistic Regression}
\eng{Logistic Regression} merupakan algoritma klasifikasi berbasis model linear yang digunakan untuk memprediksi probabilitas suatu data termasuk ke dalam kelas tertentu. Algoritma ini menggunakan fungsi sigmoid untuk mengubah keluaran linear menjadi nilai probabilitas antara 0 dan 1 \cite{ref11}. Fungsi sigmoid dinyatakan sebagai:
\begin{equation}
 \sigma(z)=\frac{1}{1+e^{-z}},
 \label{eq:sigmoid}
\end{equation}
dengan $z=\mathbf{w}^{T}\mathbf{x}+b$, di mana $\mathbf{w}$ merupakan bobot, $\mathbf{x}$ adalah fitur masukan, dan $b$ merupakan bias.

Keunggulan utama algoritma ini adalah interpretabilitas yang tinggi, karena setiap koefisien dapat menjelaskan pengaruh masing-masing fitur terhadap hasil prediksi. Namun, \eng{Logistic Regression} memiliki keterbatasan dalam menangani hubungan non-linear dan cenderung memiliki performa lebih rendah dibandingkan algoritma kompleks pada data dengan pola yang tidak sederhana \cite{ref10}.

\subsection{SVM}
\eng{Support Vector Machine} (SVM) merupakan algoritma klasifikasi yang bekerja dengan mencari \eng{hyperplane} terbaik untuk memisahkan dua kelas dengan margin maksimum. Jarak antara data terdekat dari masing-masing kelas atau \eng{support vectors} terhadap \eng{hyperplane} dibuat semaksimal mungkin \cite{ref11,ref33}. Fungsi keputusannya dapat dinyatakan sebagai:
\begin{equation}
 f(x)=\operatorname{sign}\left(\mathbf{w}\cdot\phi(x)+b\right),
 \label{eq:svm}
\end{equation}
dengan $\mathbf{w}$ adalah vektor bobot, $\phi(x)$ merupakan transformasi fitur melalui \eng{kernel}, dan $b$ adalah bias.

Keunggulan SVM terletak pada kemampuannya menangani data berdimensi tinggi dan menghasilkan margin optimal. Kelemahannya adalah waktu komputasi relatif tinggi pada dataset besar, serta kebutuhan pemilihan \eng{kernel} dan parameter yang tepat \cite{ref10}.

\subsection{KNN}
KNN merupakan algoritma berbasis \eng{instance-based learning} yang melakukan klasifikasi berdasarkan kedekatan jarak antar data. Algoritma ini menyimpan seluruh data latih dan menentukan kelas data baru berdasarkan mayoritas kelas dari tetangga terdekat \cite{ref14}. Jarak Euclidean dinyatakan sebagai:
\begin{equation}
 d(\mathbf{x},\mathbf{y})=\sqrt{\sum_{i=1}^{k}(x_i-y_i)^2}.
 \label{eq:euclidean}
\end{equation}
KNN sederhana dan mudah diimplementasikan, tetapi sensitif terhadap \eng{noise}, skala data, dan pemilihan nilai $k$. Tahap prediksi juga dapat memerlukan komputasi besar pada dataset berukuran besar \cite{ref14,ref34}.

\subsection{Gradient Boosting}
\eng{Gradient Boosting} bekerja dengan membangun model secara iteratif berdasarkan kesalahan atau residual model sebelumnya. Pada setiap iterasi, model baru dilatih untuk meminimalkan fungsi \eng{loss} dengan mengikuti arah gradien dari kesalahan yang dihasilkan \cite{ref35}. Proses tersebut dinyatakan sebagai:
\begin{equation}
 F_m(x)=F_{m-1}(x)+\gamma_m h_m(x),
 \label{eq:gradient-boosting}
\end{equation}
dengan $F_m(x)$ merupakan model pada iterasi ke-$m$, $h_m(x)$ adalah \eng{weak learner}, dan $\gamma_m$ adalah faktor pengali atau \eng{learning rate}. Algoritma ini mampu menangani hubungan non-linear, tetapi rentan \eng{overfitting} apabila parameternya tidak diatur dengan baik \cite{ref35}.

\subsection{XGBoost}
XGBoost merupakan pengembangan dari \eng{Gradient Boosting} yang dirancang untuk meningkatkan efisiensi dan performa melalui optimasi komputasi dan regularisasi \cite{ref36}. Fungsi objektifnya dapat dinyatakan sebagai:
\begin{equation}
 \mathrm{Obj}=\sum_i l(y_i,\hat{y}_i)+\sum_k\Omega(f_k),
 \label{eq:xgboost}
\end{equation}
dengan $l$ merupakan fungsi kerugian dan $\Omega$ adalah komponen regularisasi. XGBoost memiliki keunggulan dalam kecepatan pelatihan, generalisasi, dan ketahanan terhadap \eng{overfitting}, tetapi memerlukan \eng{hyperparameter tuning} yang lebih intensif dan interpretasinya lebih sulit \cite{ref36}.

\subsection{LightGBM}
LightGBM merupakan algoritma \eng{boosting} yang dikembangkan untuk meningkatkan efisiensi pelatihan, khususnya pada dataset besar. Berbeda dari pendekatan \eng{level-wise}, LightGBM menggunakan pertumbuhan pohon \eng{leaf-wise}, yaitu mengembangkan simpul yang memiliki kesalahan terbesar. LightGBM juga memanfaatkan \eng{Gradient-based One-Side Sampling} (GOSS) dan \eng{Exclusive Feature Bundling} (EFB) \cite{ref13}. Nilai \eng{gain} pada fitur ke-$j$ dengan titik \eng{split} $d$ dapat diringkas sebagai:
\begin{equation}
 V_{j\mid O}(d)=\frac{1}{n_O}\left(
 \frac{\left(\sum_{i\in O:x_{ij}\le d}g_i\right)^2}{n^{j}_{l\mid O}(d)}+
 \frac{\left(\sum_{i\in O:x_{ij}>d}g_i\right)^2}{n^{j}_{r\mid O}(d)}
 \right).
 \label{eq:lightgbm}
\end{equation}

\subsection{Optuna Framework}
\eng{Hyperparameter} merupakan parameter yang ditentukan sebelum proses pelatihan model dan sangat memengaruhi performa model \cite{ref18}. Optuna adalah kerangka optimasi \eng{hyperparameter} yang dikembangkan oleh Preferred Networks dan diperkenalkan pada 2019 \cite{ref37}. Optuna menggunakan pendekatan adaptif berbasis model probabilistik, secara bawaan menggunakan \eng{Tree-structured Parzen Estimator} (TPE), serta menyediakan mekanisme \eng{pruning} untuk menghentikan \eng{trial} yang tidak menjanjikan \cite{ref17}.

Salah satu karakteristik utama Optuna adalah prinsip \eng{define-by-run API}, sehingga ruang pencarian dapat dibangun secara dinamis. Secara umum, optimasi \eng{hyperparameter} dinyatakan sebagai:
\begin{equation}
 x^{*}=\arg\min_{x\in X} f(x),
 \label{eq:optuna}
\end{equation}
dengan $x$ adalah kombinasi parameter dan $f(x)$ merupakan fungsi objektif yang mengevaluasi performa model.

\subsection{Metrik Evaluasi}
Metrik evaluasi digunakan untuk mengukur performa model terhadap data. Dalam klasifikasi, evaluasi tidak hanya berfokus pada akurasi, tetapi juga mempertimbangkan kemampuan mendeteksi masing-masing kelas, terutama pada data tidak seimbang \cite{ref16}.

\subsubsection{Confusion Matrix}
\eng{Confusion matrix} membandingkan hasil prediksi dengan nilai aktual dan terdiri dari \eng{True Positive} (TP), \eng{True Negative} (TN), \eng{False Positive} (FP), dan \eng{False Negative} (FN) \cite{ref20}.

\begin{figure}[H]
\centering
\includegraphics[width=0.88\textwidth]{confusion_matrix.png}
\caption{\eng{Confusion Matrix} dan \eng{Performance Metrics}}
\label{fig:confusion-matrix}
\end{figure}

Metrik yang dihitung dari \eng{confusion matrix} adalah:
\begin{align}
\mathrm{Accuracy} &= \frac{TP+TN}{TP+TN+FP+FN},\\
\mathrm{Precision} &= \frac{TP}{TP+FP},\\
\mathrm{Recall} &= \frac{TP}{TP+FN},\\
F1 &= 2\times\frac{\mathrm{Precision}\times\mathrm{Recall}}{\mathrm{Precision}+\mathrm{Recall}}.
\end{align}
Pada dataset tidak seimbang, \eng{precision}, \eng{recall}, dan \eng{F1-score} dapat memberikan gambaran yang lebih realistis dibandingkan akurasi saja \cite{ref16,ref38}.

\subsubsection{Kurva ROC dan AUC}
\eng{Receiver Operating Characteristic} (ROC) menunjukkan hubungan antara \eng{True Positive Rate} dan \eng{False Positive Rate} pada berbagai nilai ambang \cite{ref21}.

\begin{figure}[H]
\centering
\includegraphics[width=0.78\textwidth]{roc_auc.png}
\caption{Kurva ROC dan AUC}
\label{fig:roc-auc}
\end{figure}

Nilai AUC yang mendekati 1 menunjukkan kemampuan klasifikasi yang semakin baik. Secara konseptual, AUC dapat dinyatakan sebagai:
\begin{equation}
 \mathrm{AUC}=\int_{0}^{1}\mathrm{TPR}\left(\mathrm{FPR}^{-1}(x)\right)\,dx.
 \label{eq:auc}
\end{equation}
Teknik \eng{cross-validation} digunakan untuk memastikan model tidak bergantung pada satu pembagian data dan menghasilkan evaluasi yang lebih stabil \cite{ref19}.

\subsection{Explainable AI}
Dalam penerapan \eng{machine learning} di bidang kesehatan, interpretabilitas model sangat penting. Model dengan akurasi tinggi tetapi sulit dipahami dapat sulit diterapkan dalam praktik klinis \cite{ref23}. \eng{Explainable AI} dikembangkan untuk meningkatkan transparansi model. Salah satu metode yang banyak digunakan adalah SHAP (SHAPley Additive exPlanations), yang menjelaskan kontribusi setiap fitur terhadap hasil prediksi \cite{ref24}. SHAP bersifat \eng{model-agnostic} dan dapat memberikan interpretasi lokal maupun global \cite{ref15}.

\begin{figure}[H]
\centering
\includegraphics[width=0.9\textwidth]{shap_visualizations.png}
\caption{Visualisasi SHAP}
\label{fig:shap}
\end{figure}

Gambar \ref{fig:shap} menunjukkan beberapa visualisasi SHAP, seperti \eng{force plot}, \eng{summary plot}, dan \eng{dependence plot}. Secara matematis, prediksi dapat dinyatakan sebagai:
\begin{equation}
 f(x)=\phi_0+\sum_{i=1}^{n}\phi_i,
 \label{eq:shap}
\end{equation}
dengan $\phi_i$ menyatakan kontribusi fitur ke-$i$ dan $\phi_0$ merupakan nilai dasar prediksi. Meskipun demikian, SHAP memiliki keterbatasan, antara lain sensitivitas terhadap distribusi data dan asumsi independensi fitur, sehingga interpretasi perlu dilakukan secara hati-hati \cite{ref25}.
